\begin{align}
\left\langle\eta_n(\mathbf r,t)\right\rangle
&\overset{\text{以分子中心为原点作逆变换(2.1.28)}}{=}
\frac{1}{(2\pi)^3}\int\mathrm d^3k\,
e^{i\mathbf k\cdot(\mathbf r-\mathbf r_n)}\tilde f(\mathbf k)\tilde\eta_n(\mathbf k,t)
\notag\\
&\overset{\text{代入 (2.1.33)}}{=}
\frac{1}{(2\pi)^3}\int\mathrm d^3k\,
e^{i\mathbf k\cdot(\mathbf r-\mathbf r_n)}\tilde f(\mathbf k)
\left\{q_n-i\mathbf k\cdot\mathbf p_n+\cdots\right\}
\tag{2.1.34}\\
&\overset{\text{拆分积分，提取 }q_n\text{、}\mathbf p_n}{=}
\frac{q_n}{(2\pi)^3}\int\mathrm d^3k\,
\tilde f(\mathbf k)e^{i\mathbf k\cdot(\mathbf r-\mathbf r_n)}
\notag\\
&\quad-\mathbf p_n\cdot\frac{1}{(2\pi)^3}\int\mathrm d^3k\,
i\mathbf k\,\tilde f(\mathbf k)e^{i\mathbf k\cdot(\mathbf r-\mathbf r_n)}+\cdots
\notag\\
&\overset{构造傅里叶逆变换}{=}
q_nf(\mathbf r-\mathbf r_n)-\mathbf p_n\cdot\nabla f(\mathbf r-\mathbf r_n)+\cdots
\tag{2.1.35}\\
&\overset{\text{(2.1.22)、(2.1.23)}}{=}
\left\langle q_n\delta(\mathbf r-\mathbf r_n)\right\rangle
-\nabla\cdot\left\langle\mathbf p_n\delta(\mathbf r-\mathbf r_n)\right\rangle+\cdots.
\tag{2.1.36}
\end{align}
